\section{An explicit exponent bound and finite weight periodicity}
\label{sec:cubic-uniform-explicit}

The finite-module conclusion of AC-K15 can be made numerical.
The only arithmetic inputs in the following argument are the
already certified base and first-layer logarithmic groups and
the generator count in AC-K16.  The remaining argument is
finite-module algebra.

\begin{theorem}[AC-K18: exponent five and nine-periodic cubic lengths]
\label{AC-K18}
For the fixed field $E=\mathbb Q(\omega,\sqrt[3]{269})$,
\[
 \lambda^5X'=0,\qquad
 \operatorname{length}_{\mathcal O}X'\leq11.           \tag{KE.1}
\]
Consequently the exponents of the individual cubic character
satisfy
\[
        3\leq m_e\leq6,\qquad m_{e+9}=m_e
                  \quad(e\geq1),                    \tag{KE.2}
\]
with $m_e=4$ for $3\mid e$.
In fact either $m_e\leq4$ for every weight, or there is one
nonzero residue class $r\pmod3$ such that
\[
 \begin{cases}
 m_e\in\{5,6\},&e\equiv r\pmod3,\\
 m_e\in\{3,4\},&e\not\equiv r\pmod3.
 \end{cases}                                         \tag{KE.3}
\]
The class $r$ and the remaining individual values are not
computed here.  For $N=t-1\geq1$, an explicit uniform budget is
\[
 3N+\lfloor N/3\rfloor
 \ \leq\sum_{e=1}^{N}m_e
 \ \leq 4N+2\lceil N/3\rceil.                       \tag{KE.4}
\]
Multiplication by $\log3$ gives the corresponding fixed
coefficient-prime-three contribution.  These statements use
no RH hypothesis.
\end{theorem}
\begin{proof}
\proofstep{1. The two conjugation components.}
Set $M=X'$, which is finite by AC-K15, and put
\[
 \delta=\omega-\omega^2,\qquad \delta^2=-3,
 \qquad B=\gamma-1,\qquad R=\mathbb Z_3[[B]].
\]
The elements $\delta$ and $\lambda$ differ by a unit.
Complex conjugation $c$ commutes with $B$ and sends
$\delta$ to $-\delta$.  The action of $B$ is nilpotent
on the finite three-group $M$: the image of the cyclotomic
pro-three group is finite, so some $\gamma^{3^h}=1$;
reduction modulo three makes $B^{3^h}=0$, and iteration
using an exponent of $M$ gives nilpotence on $M$ itself.

The quotient $M/BM$ is the degree-zero logarithmic class
group of $E$.  It has underlying group
$\mathbb Z/9\oplus\mathbb Z/3$ by AC-K13.
As a finite $\mathcal O$-module it is therefore
$\mathcal O/\delta^3$: its order is $27$, and the presence
of an element of order nine excludes every other partition
of the $\mathcal O$-length three.
Thus $M/(\delta,B)M$ is one-dimensional over $\mathbb F_3$.
Choose a lift of its generator which is a $c$-eigenvector,
by averaging with $c$ and dividing by two.  Nakayama gives
an $\mathcal O[[B]]$-generator $x$ of $M$.
Label its conjugation eigenspace by a plus sign, reversing
the two sign labels if necessary.  Then
\[
 M_+=Rx,\qquad M_-=R\delta x,\qquad M=M_+\oplus M_-.
                                                               \tag{KE.5}
\]
Indeed $\mathcal O=\mathbb Z_3\oplus\delta\mathbb Z_3$,
and averaging any coefficient with its conjugate gives
these two summands.  Multiplication by $\delta$ maps
$M_+$ onto $M_-$.
Since $(\delta^3)=(9,3\delta)$ as a $\mathbb Z_3$-lattice,
\[
             M_+/BM_+\simeq\mathbb Z/9,
       \qquad M_-/BM_-\simeq\mathbb Z/3.              \tag{KE.6}
\]
Moreover
\[
 M/\delta M\simeq M_+/3M_+,
 \qquad d:=\dim_{\mathbb F_3}(M_+/3M_+)\in\{6,7\},  \tag{KE.7}
\]
because $\delta M_+=M_-$ and
$\delta M_-=-3M_+$.  The last assertion is AC-K16.

\proofstep{2. The smaller conjugation component has order at most nine.}
Write
\[
 F(B)=(1+B)^3-1=B^3+3B^2+3B,
 \qquad C_\pm=M_\pm/F(B)M_\pm.
\]
By the first-layer descent and AC-K15,
\[
       \operatorname{length}_{\mathbb Z_3}C_+
       +\operatorname{length}_{\mathbb Z_3}C_-
                                    \leq6.           \tag{KE.8}
\]
The cyclic $R$-module $M_+/3M_+$ is
$\mathbb F_3[B]/B^d$.  As $d\geq6$ and $F(B)=B^3$
modulo three, $C_+/3C_+$ has dimension three.
Also $C_+/BC_+=M_+/BM_+\simeq\mathbb Z/9$.
Thus $C_+$ is not elementary abelian, and its length is
at least four.  Hence $C_-$ has length at most two.

Put $U=M_-$.  Its nonzero reduction modulo three is a
cyclic $\mathbb F_3[B]$-module, say
$U/3U=\mathbb F_3[B]/B^h$.
The dimension of $C_-/3C_-$ is $\min(h,3)$, so $h$ is
one or two.
If $h=2$, the length bound forces $C_-$ to be elementary
abelian of dimension two.  The surjection
$U/3U\to C_-$ is then an isomorphism, so
$F(B)U=3U$.  Since $B^2U\subset3U$,
\[
 F(B)U=B^3U+3B^2U+3BU\subset3BU.
\]
Consequently $3U=B(3U)$.  Nilpotence of $B$ forces
$3U=0$, and $U$ has order nine.

It remains to treat $h=1$.  Nakayama over $\mathbb Z_3$
makes $U$ cyclic as an abelian group: $U\simeq\mathbb Z/3^a$.
Suppose for contradiction that $a\geq3$.
The action of $B$ is multiplication by a scalar
$b\in\mathbb Z_3$.  Equation (KE.6) forces
$v_3(b)=1$.  The vector $\delta x$ generates $U$ even
as a $\mathbb Z_3$-module, and its $R$-annihilator is
$(3^a,B-b)$.
Let $H$ be the kernel of
$\delta:M_+\to U$.  It is killed by three, and as an
$\mathbb F_3[B]$-module it is generated by
\[
                    u=(B-b)x,
             \qquad z=3^a x.                          \tag{KE.9}
\]
Indeed these are the images of the generators of that
annihilator.  Also $Bz=0$:
$Bz=3^a(bx+u)=bz+3^au=0$, since $3H=0$ and $3\mid b$.
The image of $u$ in
$M_+/3M_+=\mathbb F_3[B]/B^d$ is $Bx$.
Therefore the cyclic block
$J=\mathbb F_3[B]u$ has length at least $d-1\geq5$.
Either $z\in J$, in which case it lies in the $B$-socle
of this block and hence in $B^3J$, or
$H=J\oplus\mathbb F_3z$.

Apply $F(B)$ to the exact sequence
$0\to H\to M_+\to U\to0$.
On $H$, this endomorphism is $B^3$.
Since $v_3(b)=1$, lifting the exponent gives
$F(b)=9v$ with $v\in\mathbb Z_3^\times$.
Thus $U/F(b)U$ has length two and $U[F(b)]$ is generated
by $3^{a-2}\delta x$.
The connecting map
\[
             U[F(b)]\longrightarrow H/B^3H
\]
sends this generator to the class of $vz$.
To verify the assertion, polynomial division gives
\[
 F(B)x=F(b)x+
 [B^2+Bb+b^2+3B+3b+3]u
       =F(b)x+B^2u.
\]
All terms removed in the second equality kill $u$, since
$3H=0$ and $3\mid b$.
Multiplication by $3^{a-2}$ kills the last term because
$a\geq3$, leaving $3^a vx=vz$, as claimed.

If $z\in J$, the quotient $H/B^3H$ has dimension three
and the connecting image is zero.
If $z\notin J$, that quotient has dimension four and
the connecting image has dimension one.
In either case the exact cokernel sequence gives
\[
              \operatorname{length}_{\mathbb Z_3}C_+
                         =3+2=5.
\]
But $C_-$ has length two in the present case, contradicting
(KE.8).  Hence $a\leq2$.  We have proved in both cases
\[
                          \#M_-\leq9.                \tag{KE.10}
\]

\proofstep{3. Exponent and weight consequences.}
Equation (KE.10) gives $9M_-=0$.
Since $\delta M=M_-+\delta M_-$, it follows that
$9\delta M=0$.  The identity $\delta^4=9$ now gives
$\delta^5M=0$, proving the exponent assertion in (KE.1).
Also $3M_+=-\delta M_-$ has length at most two.
By (KE.7), $M_+$ therefore has length at most $d+2$;
adding (KE.10) gives
$\operatorname{length}_{\mathcal O}M\leq d+4\leq11$.
Here the sum of the two $\mathbb Z_3$-lengths equals the
$\mathcal O$-length, since the residue field has order three.

Every $W_e(E)$ is a quotient of $M$, and AC-K12 makes it
cyclic over $\mathcal O$, of length $a_e=m_e-1$.
Thus $a_e\leq5$.  AC-K13 supplies $a_e\geq2$,
proving the bounds in (KE.2).
For $e'\equiv e\pmod9$, AC-K14 with $n=3$ identifies
$W_e/27$ and $W_{e'}/27$.  Both quotients equal the whole
group because $\lambda^5$ annihilates them and
$27$ is a unit times $\lambda^6$.  Their lengths are equal,
which proves nine-periodicity.

Finally consider the sharper alternatives.  The image of
$x$ generates $W_e$ over $\mathcal O$: after imposing
$4^e(1+B)-1=0$, the $R$-action is scalar.
Its conjugation sign is the same as that chosen in (KE.5).
Thus the minus component of the cyclic quotient
$\mathcal O/\delta^{a_e}$ has length
$\lfloor a_e/2\rfloor$.  It is also exactly
\[
                M_-/(4^e\gamma-1)M_- .              \tag{KE.11}
\]
If $M_-$ is elementary abelian, $4^e$ acts as one on it,
and (KE.6) makes (KE.11) have length one for every $e$.
Hence $a_e\in\{2,3\}$ for all $e$.
Otherwise (KE.10) forces $M_-=\mathbb Z/9$.
In this case $B$ is a scalar $b=3c\pmod9$ with
$c\not\equiv0\pmod3$, again by (KE.6).
Since $4^e\equiv1+3e\pmod9$,
\[
             4^e(1+b)-1\equiv3(e+c)\pmod9.
\]
The quotient (KE.11) has length two on the single nonzero
class $e\equiv-c\pmod3$, and length one elsewhere.
This gives (KE.3).
There are at most $\lceil N/3\rceil$ exceptional weights
among $1,\ldots,N$, proving the upper bound in (KE.4).
The lower bound uses $m_e\geq3$ and the known equality
$m_e=4$ on the $\lfloor N/3\rfloor$ multiples of three.
\end{proof}

\paragraph{What is still numerical.}
The fixed coefficient-prime-three budget is now explicit.
Determining which alternative in (KE.3) occurs requires
additional finite arithmetic data.  In either case at most
nine positive weights determine the entire sequence.
The other coefficient primes and the growing-host comparison
in H04.4 remain outside this fixed-field conclusion.
